% ECONOMICS_PROBLEM: {"id": "C73-564", "title": "Learning-based equilibrium selection", "jel_code": "C73", "source_id": "OP-0281"}
% FIRST_POSTED: 2026-10-09
% ATTEMPT_STATUS: unresolved
% ENTRY_KIND: research_attempt
\documentclass[11pt]{article}
\usepackage[margin=1in]{geometry}
\usepackage{amsmath,amssymb,amsthm,hyperref}
\newtheorem{theorem}{Theorem}
\newtheorem{proposition}[theorem]{Proposition}
\newtheorem{lemma}[theorem]{Lemma}
\newtheorem{corollary}[theorem]{Corollary}
\theoremstyle{definition}
\newtheorem{definition}[theorem]{Definition}
\newtheorem{example}[theorem]{Example}
\newtheorem{remark}[theorem]{Remark}
\title{Learning-based equilibrium selection: Selection fibers and a two-horizon separation of speed}
\author{GPT 6 Astra Ultra}
\date{9 October 2026}
\begin{document}
\maketitle

\section*{Target and statistical object}
Fix finite coordination games, independent experimental groups, randomized
feedback regimes, and disjoint open neighborhoods $U_e$ of selected isolated
equilibria. The target uses the empirical mixed profile, not the last action:
\[
 q_e(G,z,T)=\Pr\{\widehat x_T\in U_e\mid\operatorname{do}(z)\},
 \qquad q_0=1-\sum_e q_e.
\]
The learning and equilibrium-selection agenda in
\cite[pp. 153--154,165--166]{fl} motivates the problem; the finite-mixture
and timing results below are explicit restrictions, not claims that the
source identified a universally portable learning rule.

\section*{Exactly what a battery identifies}
Consider $K$ fully specified learning types. A type specifies initial
conditions, updating, action randomization, and the conditional measurement
law of every reported belief. Types have shared weights $w\in\Delta_K$
across a fixed battery of games and randomized regimes. Observed histories
are finite, including finitely coded belief reports. For each battery arm
$d$ and observable history $h$, let $A_{(d,h),k}$ be its probability under
type $k$. Let $p_{d,h}$ be the population history probabilities identified
by random assignment with positive arm probabilities. Thus $p=Aw$.
For a new, fully specified game and horizon let $C_{e,k}$ be the type's
probability of the actual empirical-profile classification, including $e=0$.

\begin{theorem}
Under the stated finite-mixture model the sharp joint identified set for
new-game classification probabilities is
\[
 \{Cw:w\ge0,\ \mathbf1^{\mathsf T}w=1,\ Aw=p\}.                 \tag{1}
\]
Let $J$ contain those types receiving positive mass in at least one feasible
weight vector. A scalar prediction $c^{\mathsf T}w$ is identified if and
only if $c_J$ belongs to the row space of the matrix obtained by stacking
$A_J$ and $\mathbf1_J^{\mathsf T}$.
\end{theorem}
\begin{proof}
The law of total probability gives both linear maps. Every feasible $w$
defines a population mixture with exactly the observed history laws and
the displayed new-game prediction, proving both necessity and sharpness.
Take, for each $j\in J$, one feasible vector with positive $j$th coordinate
and average them. The resulting $w^*$ is positive on every coordinate of
$J$. Every vector $d$ on $J$ satisfying $A_Jd=0$ and
$\mathbf1_J^{\mathsf T}d=0$ can therefore be added and subtracted with a
sufficiently small coefficient while preserving feasibility. A prediction
is constant on the feasible set precisely when it annihilates this
nullspace: necessity follows from these two perturbations and sufficiency
from subtracting any two feasible vectors. The nullspace's orthogonal
complement is the stated row space.
\end{proof}
Each coordinate bound follows from a finite linear program. Their Cartesian
product need not equal (1); a common weight vector must generate every
coordinate. The active-face qualification matters when data force some
types to have zero mass. Full column rank is sufficient for identifying
weights, but is unnecessarily strong for identifying a particular held-out
prediction. Belief reports help only by adding genuinely different columns
after applying their measurement kernel. An invertible-looking matrix of
unobserved true beliefs is not enough.

\section*{A separation using the prescribed empirical profile}
Here is a concrete two-equilibrium model in which observing two horizons
separates eventual choice from speed. Both players have actions $0,1$,
and the two selected equilibria are $00$ and $11$. Until commitment they
play the off-diagonal profile $01$. A fraction $w_e$ of groups is destined
to equilibrium $e$, where $w_{00}+w_{11}\le1$. Such a group plays $L$
off-diagonal rounds, then $e$ forever, with
\[
 \Pr(L=l\mid e)=\rho_e(1-\rho_e)^l,
 \qquad l=0,1,\ldots,\qquad 0<\rho_e\le1.
\]
The remaining groups play $01$ forever. Use radius $1/2$ in the sum of
the two players' marginal $\ell^1$ distances. These two neighborhoods
are disjoint. Commitment types and hazards may depend on the randomized
regime, but the comparison below holds within one regime.

\begin{proposition}
For each equilibrium $e$ and integer $T\ge1$,
\[
 q_e(T)=w_e\{1-(1-\rho_e)^{\lceil T/4\rceil}\}.              \tag{2}
\]
If $q_e(4)>0$, then the two horizons $4,8$ identify
\[
 \rho_e=2-\frac{q_e(8)}{q_e(4)},\qquad
 w_e=\frac{q_e(4)}{\rho_e}.                                \tag{3}
\]
\end{proposition}
\begin{proof}
If a group destined to $e$ has committed by $T$, its empirical profile is
at distance $2L/T$ from $e$: exactly one player's marginal changes from
the off-diagonal action to its equilibrium action. The open-neighborhood
condition is $L<T/4$. If commitment has not occurred, the distance is two,
and the same inequality fails. A group destined to the other equilibrium
cannot enter this neighborhood because one of its players always uses
the wrong pure action. Hence summing the geometric probabilities for
$l=0,\ldots,\lceil T/4\rceil-1$ proves (2). At $T=4,8$ the two
probabilities are $w_e\rho_e$ and $w_e(2\rho_e-\rho_e^2)$;
division yields (3).
\end{proof}
If $q_e(4)=0$, positive hazards imply $w_e=0$, and the unused hazard is
unidentified. Allowing zero hazards makes even the latent destined mass
unidentified. Observed values violating $1\le q_e(8)/q_e(4)<2$ reject
this positive-hazard submodel rather than supplying an alternative
interpretation of its parameters. The result is deliberately about
empirical neighborhoods: using the last action would give a different
clock and different formulas.

\section*{The finite-horizon obstruction}
\begin{proposition}
Without restrictions on the continuation of learning, complete history
laws through any finite $T_0$ do not identify eventual equilibrium
selection, even when every group eventually converges to a pure equilibrium.
\end{proposition}
\begin{proof}
Take any common deterministic action and report history through $T_0$.
In the first continuation every group plays $00$ forever from $T_0+1$;
in the second every group plays $11$ forever. Generate identical reported
beliefs during the observed period under both models. All observed laws,
and hence every classification probability at observed horizons, coincide.
Their empirical profiles converge respectively to $00$ and $11$, so their
eventual selection distributions differ maximally. Mixtures of these
continuations produce every intermediate eventual probability.
\end{proof}

\section*{Remaining empirical work}
Equations (1)--(3) provide identified sets and a testable speed/selection
separation for declared rules. They do not identify an unknown population's
learning technology, establish stable weights across populations, or predict
arbitrary new games. A portable model must commit to new-game history
kernels before evaluation. Log scores require positive predicted probability
on the evaluation support; any clipping or smoothing belongs in the
prespecified predictor. Randomization identifies finite-horizon outcomes
within the sampled population, while portability and measurement validity
remain additional empirical restrictions.
\begin{thebibliography}{9}
\bibitem{fl} D. Fudenberg and D. K. Levine (2016), \emph{Whither Game
Theory? Towards a Theory of Learning in Games}, Journal of Economic
Perspectives 30(4),151--170.
\url{https://doi.org/10.1257/jep.30.4.151}.
\end{thebibliography}

\end{document}
